\title{FlashAttention-3: Fast and Accurate Attention with Asynchrony and Low-precision}

\begin{abstract}
Within the Transformer framework, attention mechanisms constitute a pivotal yet performance-limiting component, especially for large language models and tasks requiring extensive context windows. While \textsc{FlashAttention} proposed a method to accelerate attention computation on GPUs by reducing memory traffic, its design has yet to fully harness advancements in contemporary GPU hardware. For instance, \textsc{FlashAttention-2} only attains a 35\% utilization rate on the H100 GPU. In this work, we introduce three primary strategies to further optimize attention on Hopper GPUs: leveraging asynchronous features of Tensor Cores and TMA to (1) concurrently conduct computation and data transfer using warp specialization, (2) interleave matrix multiplication and softmax computation at the block level, and (3) utilize block-wise quantization and non-uniform processing enabled by hardware-based support for FP8 precision. Our approach, termed \textsc{FlashAttention-3}, achieves between 1.5 and 2.0$\times$ speedup on H100 GPUs; with FP16, performance peaks at 740 TFLOPs/s (a 75\% utilization rate), while FP8 approaches 1.2 PFLOPs/s. We also show that the FP8 version of \textsc{FlashAttention-3} produces 2.6$\times$ less numerical error compared to a standard FP8 attention implementation.
\end{abstract}

\section{Introduction}
\label{sec:intro}

Within the Transformer architecture~\citep{vaswani2017attention}, the primary computational constraint arises from the attention mechanism, as calculating the self-attention scores between queries and keys increases quadratically with respect to sequence length. Expanding attention to handle longer contexts stands to facilitate novel capabilities, such as enabling models to process and reason across several lengthy documents~\citep{guo2021longt5,shaham2022scrolls,peng2023yarn}, or manage multiple files in extensive codebases~\citep{roziere2023code, li2023starcoder}. Moreover, it opens possibilities for new data modalities, including high-resolution images~\citep{chen2022scaling}, audio~\citep{gulati2020conformer}, and video~\citep{ho2022video}, as well as emerging applications, like engaging with users with substantial conversational histories~\citep{sun2019bert4rec} or supporting agents through extended task sequences~\citep{yao2022react}. Consequently, there has been heightened focus on enhancing the efficiency of attention for long contexts—through various approximation approaches~\citep{katharopoulos2020transformers,choromanski2020rethinking,
tay2020efficient}, software-level acceleration (\citep{rabe2021self,dao2022flashattention,kwon2023efficient}), and by devising alternative neural architectures~\citep{peng2023rwkv,sun2023retentive,gu2023mamba}.

This study extends the research by \citet{dao2022flashattention}, who advanced the development of exact-attention algorithms by embedding an awareness of GPU execution paradigms and hardware details into algorithmic design. \citet{dao2022flashattention} proposed \textsc{FlashAttention}, which leverages a novel tiling technique to parallelize the attention mechanism while circumventing costly intermediate accesses to global memory. This is achieved by integrating the entire set of attention computations into a unified GPU kernel. Subsequently, \citet{dao2023flashattention2} enhanced the algorithm, releasing \textsc{FlashAttention-2}, which incorporates parallelization along the sequence length and organizes the key and value matrix computations in the forward pass into block-wise processing. This adaptation aims to maximize GPU workload distribution and occupancy. Despite these advances, our evaluation indicates that \textsc{FlashAttention-2} still demonstrates significantly lower hardware utilization on modern GPUs in comparison with optimized GEMM operations—achieving only 35\% utilization on the Hopper H100, whereas GEMM kernels reach 80-90\%. Some of this inefficiency arises from implementation choices; for example, \textsc{FlashAttention-2} does not fully exploit Hopper-specific instructions for Tensor Cores, instead relying on instructions tailored for the previous Ampere architecture. Recent projects like ThunderKitten~\citep{spector2024thunder} and cuDNN 9~\citep{cudnn9} illustrate that, by leveraging instructions unique to Hopper and adopting tile-based algorithmic approaches, attention computation can be both accelerated and simplified.

At a more foundational level, the algorithm underlying \textsc{FlashAttention-2} is structured around a conventional synchronous paradigm, without explicitly exploiting asynchronous execution or low-precision arithmetic within its architecture. Asynchrony emerges due to the presence of hardware components specifically optimized to expedite key operations within machine learning tasks: specialized units such as Tensor Cores handle matrix multiplications, and the Tensor Memory Accelerator (TMA) manages memory loads—distinct from standard CUDA cores that execute logical, integer, or floating-point operations. Low-precision formats—like FP8 on Hopper and FP4 on Blackwell—continue the lineage established by FP16 (introduced with Pascal in 2017) and BF16 (introduced on Ampere in 2020), demonstrating substantial increases in computational throughput (by two or four times) while maintaining the same chip area and power demands. Section \cref{subsec:hardware} surveys the advancements that Hopper introduces in these areas. The principal technical obstacle is revising \textsc{FlashAttention-2} so it can take full advantage of such hardware properties: realizing asynchrony involves concurrently executing matmul and softmax (despite their data dependencies), while implementing low-precision methods necessitates careful error management to suppress quantization artifacts, particularly given the prevalence of outlier features in LLMs~\citep{dettmers2208llm, sun2024massive}.

To address these challenges, we introduce \textsc{FlashAttention-3}, an approach that unifies three novel techniques aimed at optimizing performance on state-of-the-art GPU platforms.\footnote{Although our experimental context focuses on NVIDIA's Hopper architecture, the underlying algorithm is applicable to any GPU system that supports robust asynchronous operations and low-precision computations.}

\begin{enumerate}

\item \textbf{Producer-Consumer asynchrony:} We introduce a warp-specialized software pipelining approach that takes advantage of the asynchronous nature of data movement and Tensor Core computation. By allocating data producers and consumers to distinct warps, this scheme increases the opportunity for the algorithm to mask latencies stemming from memory access and instruction issuance.
\item \textbf{Hiding softmax under asynchronous block-wise GEMMs:} Our method interleaves the low-throughput non-GEMM operations used in softmax (such as floating-point multiply-adds and exponentials) with the execution of asynchronous WGMMA instructions for GEMM tasks. To facilitate this, we restructure the \textsc{FlashAttention-2} algorithm to break select sequential dependencies linking the softmax and GEMM stages. For instance, in the two-stage algorithmic variant, as the softmax runs on one scores matrix block, the WGMMA performs asynchronously to prepare computation for the subsequent block.
\item \textbf{Hardware-accelerated low-precision GEMM:} We modify the forward pass strategy so it can utilize FP8 Tensor Cores for GEMM operations, which leads to an observed near doubling of TFLOPs/s throughput. This adaptation necessitates reconciling WGMMA’s expectations for the memory layout of FP32 accumulator and FP8 operand blocks. Techniques such as block quantization and incoherent processing are employed to alleviate the accuracy degradation associated with computation in FP8 format.
\end{enumerate}

To empirically substantiate our approach, we conducted benchmarks of \textsc{FlashAttention-3} on the H100 SXM5 GPU across diverse parameter settings. The results indicate that (1) with FP16 precision, the forward pass achieves a 1.5-2.0$\times$ speedup compared to \textsc{FlashAttention-2} (reaching peak throughput of 740 TFLOPs/s) and a 1.5-1.75$\times$ increase during the backward pass; (2) FP8 precision attains performance near 1.2 PFLOPs/s; and (3) when sequence lengths are large, FP16 shows superior performance while FP8 remains competitive\footnote{Specifically, with head dimension 64, \textsc{FlashAttention-3} FP8 demonstrates better performance, whereas with head dimensions of 128 and 256 it is on par in scenarios without causal masking and lags when causal masking is applied.} relative to the current state-of-the-art attention kernel from NVIDIA’s cuDNN library. We further confirmed that FP16 \textsc{FlashAttention-3} produces numerical errors comparable to those of \textsc{FlashAttention-2} and surpasses the canonical attention implementation, as intermediate computations such as softmax normalization remain in FP32 format. Additionally, under outlier feature conditions, FP8 \textsc{FlashAttention-3} employing block quantization and incoherent processing achieves 2.6$\times$ greater accuracy than conventional attention using per-tensor quantization.

We have made \textsc{FlashAttention-3} publicly available under a permissive open-source license\footnote{\textsc{FlashAttention-3} is accessible at \texttt{https://github.com/Dao-AILab/flash-attention}} and intend to integrate it with the PyTorch and Hugging Face ecosystems in order to maximize its accessibility for both researchers and practitioners.

\section{Background: Multi-Head Attention and GPU Characteristics}
\label{sec:background}

\subsection{Multi-Head Attention}
\label{subsec:multi_head_attn}

Consider $\mathbf{Q}, \mathbf{K}, \mathbf{V} \in \mathbb{R}^{N \times d}$, representing the query, key, and value matrices for a single attention head, where $N$ denotes the sequence length and $d$ the dimensionality of each head. The output of the attention mechanism, $\mathbf{O}$, is obtained as follows:

\begin{equation*}
  \mathbf{S} = \alpha \mathbf{Q} \mathbf{K}^\top \in \mathbb{R}^{N \times N}, \quad \mathbf{P} = \mathrm{softmax}(\mathbf{S}) \in \mathbb{R}^{N \times N}, \quad \mathbf{O} = \mathbf{P}\mathbf{V} \in \mathbb{R}^{N \times d},
\end{equation*}

Here, the $\mathrm{softmax}$ function is computed across each row, and the scaling factor $\alpha$ is usually chosen as $1/\sqrt{d}$. To avoid numerical issues related to large exponentials, it is common to subtract $\mathrm{rowmax}(\mathbf{S})$ from $\mathbf{S}$ before applying the softmax. In the case of multi-head attention (MHA), independent projections for queries, keys, and values are used for each head, enabling this process to be executed concurrently over different heads and batches to yield the complete output tensor.

Consider a scalar loss function $\phi$, and denote the gradient with respect to a quantity $(-)$ as $\mathbf{d}(-) = \partial \phi / \partial (-)$. Given the gradient of the output, $\mathbf{dO} \in \mathbb{R}^{N \times d}$, we derive $\mathbf{dQ}$, $\mathbf{dK}$, and $\mathbf{dV}$ by systematically applying the chain rule:

\begin{align*}
  \mathbf{dV} &= \mathbf{P}^\top \mathbf{dO} \in \mathbb{R}^{N \times d} \\
  \mathbf{dP} &= \mathbf{dO} \mathbf{V}^\top \in \mathbb{R}^{N \times N} \\
  \mathbf{dS} &= \mathrm{dsoftmax} (\mathbf{dP}) \in \mathbb{R}^{N \times N} \\
  \mathbf{dQ} &= \alpha \mathbf{dS} \mathbf{K} \in \mathbb{R}^{N \times d} \\
  \mathbf{dK} &= \alpha \mathbf{dS}^\top \mathbf{Q} \in \mathbb{R}^{N \times d},
\end{align*}

For a vector $s$ and $p = \mathrm{softmax}(s)$, the derivative is computed as $\mathbf{d}s = (\mathrm{diag}(p) - p p^\top)\mathbf{d}p$, which is then applied across each row when denoted $\mathrm{dsoftmax}(\mathbf{dP})$. Notably, these computations are amenable to parallelization over the different heads and batches during the backward phase of MHA.

\subsection{GPU hardware characteristics and execution model}
\label{subsec:hardware}

We describe the aspects of the GPU's execution model relevant for \textsc{FlashAttention-3}, with a focus on the NVIDIA Hopper architecture as a concrete instantiation of this model.

\paragraph{Memory hierarchy:} The memory architecture of the GPU is structured as a hierarchical system of data storage regions, where available capacity is inversely proportional to the memory bandwidth (\cref{tab:gpu-hierarchy})\footnote{\citet{luo2024benchmarking} indicates that shared memory achieves a bandwidth of 128 bytes per clock cycle per SM, with the total estimate calculated based on 132 SMs and a boost clock rate of 1830 MHz.}. The lowest level of the hierarchy is global memory (GMEM), also referred to as HBM, which constitutes off-chip DRAM accessible by all streaming multiprocessors (SMs). Data retrieved from GMEM is automatically cached within an on-chip L2 cache layer. Beyond this, each SM is equipped with a dedicated, high-banking, on-chip cache known as shared memory (SMEM), which is directly managed by the programmer. The hierarchy culminates at the register file, located within each SM.

\paragraph{Thread hierarchy:} The GPU organizes its computation using a structured hierarchy of threads. At the most granular level are individual threads, which are grouped into warps, each containing 32 threads. Four consecutive warps constitute a warpgoup. These warpgroups are further organized into threadblocks, also referred to as cooperative thread arrays (CTAs). In architectures such as Hopper, threadblocks can be aggregated into threadblock clusters, and at the highest level, all these elements are assembled into grids.

The two hierarchies exhibit a strong interconnection. Threads that belong to the same CTA are scheduled together on a single SM, while CTAs within one cluster are executed together on a single GPC. All threads inside a CTA have direct access to SMEM, while individual threads are limited to a maximum of 256 private registers (RMEM) each.

\paragraph{Asynchrony and warp-specialization:}

GPUs function as throughput-oriented processors, leveraging concurrent execution and asynchronous operations to mask delays arising from memory accesses and computation. In the case of asynchronous memory transfers between GMEM and SMEM, Hopper introduces the Tensor Memory Accelerator (TMA), a specialized hardware component \cite[\S7.29]{cuda}. Additionally, departing from earlier architectures like Ampere, Hopper's Tensor Core—accessible through the warpgroup-wide WGMMA instruction \cite[\S9.7.14]{ptx}—operates asynchronously as well, enabling it to fetch input data directly from shared memory.

Hardware mechanisms that support asynchrony enable warp-specialized kernels, wherein the warps within a CTA are assigned exclusive roles as either producers or consumers, resulting in each warp being dedicated to either computation or data transfer tasks exclusively. This approach generally enhances the compiler's capacity to emit more efficient instruction schedules \citep{warp-specialization-2011}. Furthermore, Hopper introduces functionality to dynamically redistribute register resources among warpgroups using \verb|setmaxnreg| \cite[\S9.7.17.1]{ptx}, thereby allowing warps engaged in MMA operations to access a greater portion of RMEM compared to warps responsible solely for issuing TMA instructions (which can be managed by just a single thread).

\paragraph{Low-precision number formats:}
\label{sec:low-precision-gpu}
Contemporary GPUs are equipped with dedicated hardware to accelerate computations using low-precision numerical formats. For instance, on Hopper, the WGMMA instruction can utilize FP8 Tensor Cores, achieving twice the per-SM throughput relative to operations performed using FP16 or BF16 formats.

Nonetheless, utilizing FP8 WGMMA correctly requires a detailed understanding of its operand layout requirements. Consider the case of performing a GEMM operation to compute $A \times B^{\top}$, where $A$ is an $M \times K$ matrix and $B$ is an $N \times K$ matrix. We define an operand as \emph{mn-major} if it is stored contiguously along its outer $M$ or $N$ dimension, and as \emph{k-major} if the contiguous dimension is the inner $K$ dimension. With FP16 WGMMA, input operands held in SMEM may be supplied in either mn-major or k-major layouts. By contrast, FP8 WGMMA imposes a stricter requirement, allowing solely k-major inputs. This restriction introduces complications, for example in scenarios such as attention mechanisms, where one might seek to combine multiple GEMMs within a single kernel; here, incompatible memory layouts between FP32 accumulators and FP8 operands create challenges in launching subsequent dependent FP8 WGMMAs.

In the context of attention, these layout restrictions entail certain modifications to the design of an FP8 algorithm, which we describe in \cref{sec:algofp8}.

\subsection{Standard Attention and Flash Attention} In line with \citet{dao2022flashattention}, we refer to \textbf{standard attention} as the GPU-based implementation of attention where the intermediate matrices $\mathbf{S}$ and $\mathbf{P}$ are stored in HBM. The primary innovation introduced by \textsc{FlashAttention} involves utilizing a local approach to the softmax reduction, thereby minimizing costly intermediate memory operations and merging the attention computation within a single GPU kernel. This local softmax procedure is represented by lines \ref{code-ws:softmax_start}-\ref{code-ws:softmax_end} in the consumer mainloop of \cref{alg:flash3_wgmma_ws_only}, along with the necessary block-wise rescaling of $\mathbf{O}$. A concise explanation demonstrating that this approach faithfully computes $\mathbf{O}$ is provided in \cite[\S 2.3.1]{dao2023flashattention2}.

\section{FlashAttention-3: Algorithm}
\label{sec:algo}

This section outlines the \textsc{FlashAttention-3} algorithm. We concentrate on the forward pass for clarity; details regarding the backward pass algorithm can be found in~\cref{sec:algo_ws_bwd}. Initially, we describe the integration of warp-specialization with a circular SMEM buffer into the foundational framework of \textsc{FlashAttention-2}. Next, we discuss leveraging the asynchrony of WGMMA to construct a two-stage GEMM-softmax pipeline with overlapping stages. Finally, we elaborate on the adjustments necessary for FP8 support, encompassing both the enforcement of appropriate layout and improvements in accuracy using block quantization and incoherent processing.

\subsection{Producer-Consumer asynchrony through warp-specialization and pingpong scheduling}
\label{sec:algo_ws}

\paragraph{Warp-specialization}
Similar to \textsc{FlashAttention-2}, \textsc{FlashAttention-3} achieves extensive parallelization across the dimensions of batch size, number of attention heads, and length of the query sequence during the forward pass. Therefore, it is sufficient to analyze the algorithm from the perspective of a CTA, which processes an individual tile $\mathbf{Q}_i$ of the query matrix in order to produce the respective output tile $\mathbf{O}_i$. For clarity in the exposition, we initially describe the warp-specialization strategy using a circular SMEM buffer, without introducing additional GEMM-softmax overlap. Let the head dimension be denoted by $d$ and the sequence length by $N$. By fixing the query block size as $B_r$, the query matrix $\mathbf{Q}$ is partitioned into $T_r = \lceil \frac{N}{B_r} \rceil$ blocks labeled $\mathbf{Q}_1, .., \mathbf{Q}_{T_r}$.

\begin{algorithm}[H]
    \caption{\small\label{alg:flash3_wgmma_ws_only}\textsc{FlashAttention-3} forward pass \textbf{excluding} intra-consumer overlap -- CTA-oriented perspective}
    \begin{algorithmic}[1]
\REQUIRE Matrices $\mathbf{Q}_i \in \mathbb{R}^{B_r \times d}$ and $\mathbf{K}, \mathbf{V} \in \mathbb{R}^{N \times d}$ resident in HBM; specify key block size $B_c$, yielding $T_c = \lceil \frac{N}{B_c} \rceil$.
\STATE Construct a pipeline object that coordinates barrier synchronization over an $s$-stage circular buffer allocated in shared memory.
\IF {executing within the producer warpgroup}
\STATE Free a predefined quota of registers.
\STATE Trigger data transfer for $\mathbf{Q}_i$ from HBM to shared memory.
\STATE Once transfer completes, signal consumers that $\mathbf{Q}_i$ is available.
\FOR{$0 \le j < T_c$}
    \STATE Await consumption of the $(j\,\%\,s)$th buffer stage.
    \STATE Initiate transfer of $\mathbf{K}_j, \mathbf{V}_j$ from HBM to the $(j\,\%\,s)$th position in shared memory.
    \STATE Upon completion, notify consumer warps of $\mathbf{K}_j, \mathbf{V}_j$ availability.
\ENDFOR
\ELSE \STATE Reclaim the reserved registers according to consumer warp count.
\STATE Allocate and initialize on-chip data: $\mathbf{O}_i = (0) \in \mathbb{R}^{B_r \times d}$ along with $\ell_i = (0)$, $m_i = (-\infty)$ in $\mathbb{R}^{B_r}$.
\STATE Wait until $\mathbf{Q}_i$ is available in shared memory.
\FOR{$0 \le j < T_c$}
\STATE Wait for the completion of $\mathbf{K}_j$ transfer into shared memory.
\STATE Execute $\mathbf{S}_i^{(j)} = \mathbf{Q}_i \mathbf{K}_j^T$ using SS-GEMM, commit results and synchronize. \label{alg:ws_only_gemm1}
\STATE Preserve previous $m_i$ in $m_i^{\mathrm{old}}$; update $m_i$ via $m_i = \mathrm{max}(m_i^{\mathrm{old}}, \mathrm{rowmax}(\mathbf{S}_i^{(j)}))$. \label{code-ws:softmax_start}
\STATE Evaluate $\widetilde{\mathbf{P}}_i^{(j)} = \mathrm{exp}(\mathbf{S}_i^{(j)} - m_i)$; refresh $\ell_i$ as $\ell_i = \mathrm{exp}(m_i^{\mathrm{old}} - m_i)\,\ell_i + \mathrm{rowsum}(\widetilde{\mathbf{P}}_i^{(j)})$. \label{code-ws:softmax_end}
\STATE Wait for $\mathbf{V}_j$ to appear in shared memory.
\STATE Perform the operation $\mathbf{O}_i = \mathrm{diag}(\mathrm{exp}(m_i^{\mathrm{old}} - m_i))^{-1}\mathbf{O}_i + \widetilde{\mathbf{P}}_i^{(j)} \mathbf{V}_j$ using RS-GEMM, then commit and synchronize. \label{alg:ws_only_gemm2}
\STATE Hand control of the $(j\,\%\,s)$th buffer stage back to the producer.
\ENDFOR
\STATE Compute outputs via $\mathbf{O}_i = \mathrm{diag}(\ell_i)^{-1} \mathbf{O}_i$ and $L_i = m_i + \log(\ell_i)$.
\STATE Write $\mathbf{O}_i$ and $L_i$ to HBM, storing them as the $i$th partitions of $\mathbf{O}$ and $L$.
\ENDIF
\end{algorithmic}
\end{algorithm}

In our implementation of \cref{alg:flash3_wgmma_ws_only} targeting Hopper, we employ \verb|setmaxnreg| for allocating and deallocating resources, utilize TMA to load $\mathbf{Q}_i$ as well as $\{ \mathbf{K}_j, \mathbf{V}_j \}_{0 \leq j < T_c}$, and apply WGMMA for carrying out GEMMs within the consumer mainloop. Here, the SS or RS prefix designates whether shared memory or register file supplies the first operand, respectively. It is important to understand, when analyzing the control flow of \cref{alg:flash3_wgmma_ws_only}, that TMA loads operate asynchronously, so launching a TMA load does not block on previous ones finishing. Additionally, during the first $s$ iterations of the producer mainloop, waits are not introduced, since these steps are necessary for initially populating the buffer.

\paragraph{Pingpong scheduling} The asynchronous characteristics of WGMMA and TMA, combined with warp specialization, allow for concurrent execution of the softmax operation in one warpgroup while another warpgroup performs the GEMM operation. To illustrate the need for such overlapping, observe that on contemporary hardware accelerators, the performance of non-matrix multiplication kernels is significantly lower than that of matmul kernels. For instance, on the H100 SXM5 GPU, FP16 matmul achieves a peak throughput of 989 TFLOPS, whereas special function units (responsible for computations like exponentials\footnote{According to the CUDA programming guide, each streaming multiprocessor (SM) can carry out 16 special function operations per clock cycle. Multiplying 16 operations by 132 SMs and a 1830 MHz clock speed results in a peak rate of 3.9 TFLOPS for special functions.}) only reach 3.9 TFLOPS. During the attention forward computation in FP16 for a head size of 128, the workload consists of 512 times more FLOPs dedicated to matmul compared to those involving exponentials, but the execution rate of exponentials is 256 times slower than matmul, so computing the exponentials may occupy up to half the execution time of a single cycle. The imbalance becomes even more pronounced in FP8 precision: although the throughput for matmul doubles, the rate at which exponentials can be computed does not improve.

Given that the exponential computation is delegated to a dedicated hardware component (the multi-function unit), it is preferable to overlap its execution with the matrix multiplication operations carried out by the Tensor Cores. To orchestrate this, synchronization barriers (\texttt{bar.sync} instructions) are utilized, ensuring that the GEMM operations (specifically, GEMM1 -- $\mathbf{P} \mathbf{V}$ from the current iteration, and GEMM0 -- $\mathbf{Q} \mathbf{K}^\top$ from the subsequent iteration) associated with warpgroup 1 are scheduled ahead of those belonging to warpgroup 2. This arrangement allows the softmax phase for warpgroup 1 to coincide temporally with the GEMMs for warpgroup 2. Subsequently, the two warpgroups exchange functions, so that warpgroup 2 proceeds with softmax while warpgroup 1 resumes GEMM computation—a process referred to as ``pingpong'' scheduling, as depicted in~\cref{fig:pingpong_scheduling}. Although the pingpong scheduling mechanism does not always execute in the exact alternating fashion suggested in the diagram, this methodology typically yields substantial performance gains (e.g., elevating throughput from 570 TFLOPS to between 620 and 640 TFLOPS for FP16 forward passes, with a head size of 128 and sequence length set at 8192).

\paragraph{Attention variants} In the cases of multi-query attention~\citep{shazeer2019fast} and grouped query attention~\citep{ainslie2023gqa}, the methodology implemented in \textsc{FlashAttention-2} is followed, adapting the tensor indexing strategy accordingly to circumvent unnecessary replication of $\mathbf{K}$ and $\mathbf{V}$ within HBM.

\subsection{Intra-warpgroup overlapping GEMMs and softmax}

Even within one warpgroup, we can overlap some instructions in the softmax with some instructions in the GEMMs. We describe one technique to do so.

Within the attention algorithm, there are inherent sequential dependencies among operations in the inner loop (main loop), limiting intra-iteration parallelization. As an illustration, computing (local) softmax (lines \ref{code-ws:softmax_start} to \ref{code-ws:softmax_end}) depends on the result $\mathbf{S}_i^{(j)}$ from the initial GEMM operation, and subsequently, the second GEMM requires the outcome $\widetilde{\mathbf{P}}_i^{(j)}$ from the softmax step. The wait instructions found on lines \ref{alg:ws_only_gemm1} and \ref{alg:ws_only_gemm2} in \cref{alg:flash3_wgmma_ws_only} ensure that both the softmax and GEMM procedures execute in a strictly ordered sequence. Nevertheless, these inter-stage dependencies can be circumvented by introducing pipelining across loop iterations, utilizing extra register-based buffers. Building on this concept, we introduce a novel two-stage\footnote{It is important to recognize that while the number of pipeline stages is constrained by the circular SMEM buffer’s stage count $s$, it does not necessarily have to match it exactly.} GEMM-softmax pipelining scheme:

\begin{algorithm}[H]
    \caption{\small\label{alg:flash3_wgmma}\textsc{FlashAttention-3} consumer warpgroup forward pass}
    \begin{algorithmic}[1]
      \REQUIRE  Matrices $\mathbf{Q}_i \in \mathbb{R}^{B_r \times d}$ along with $\mathbf{K}, \mathbf{V} \in \mathbb{R}^{N \times d}$ residing in HBM, a key block size $B_c$, and $T_c = \lceil \frac{N}{B_c} \rceil$.
      \STATE Dynamically allocate registers according to the active number of consumer warps.
      \STATE Within on-chip memory, set $\mathbf{O}_i = (0) \in \mathbb{R}^{B_r \times d}$, while initializing $\ell_i$ and $m_i$ as $(0)$ and $(-\infty)$ respectively in $\mathbb{R}^{B_r}$.
      \STATE Await the loading of $\mathbf{Q}_i$ and $\mathbf{K}_0$ into shared memory.
      \STATE Utilize WGMMA to calculate $\mathbf{S}_{\mathrm{cur}} = \mathbf{Q}_i \mathbf{K}_0^T$, ensuring completion by waiting. 
      \STATE Free up the buffer at stage $0$ associated with $\mathbf{K}$.
      \STATE Derive $m_{i}$, $\tilde{\mathbf{P}}_{\mathrm{cur}}$, and $\ell_{i}$ from $\mathbf{S}_{\mathrm{cur}}$, then appropriately scale $\mathbf{O}_i$.
      \FOR{$1 \le j < T_c - 1$}
        \STATE Pause execution until $\mathbf{K}_j$ arrives in shared memory.
        \label{code:mainloop_start}
        \STATE Execute WGMMA for $\mathbf{S}_{\mathrm{next}} = \mathbf{Q}_i \mathbf{K}_{j}^T$; submit the computation but do not block.
        \label{code:first_wgmma}
        \STATE Hold until $\mathbf{V}_{j-1}$ is accessible in shared memory.
        \STATE Use WGMMA to perform $\mathbf{O}_{i} = \mathbf{O}_{i} + \tilde{\mathbf{P}}_{\mathrm{cur}} \mathbf{V}_{j-1}$, issuing but not synchronizing.
        \label{code:second_wgmma}
        \STATE Wait for the $\mathbf{Q}_i \mathbf{K}_{j}^T$ WGMMA computation to finalize.
        \STATE Update $m_{i}$, $\tilde{\mathbf{P}}_{\mathrm{next}}$, and $\ell_{i}$ by processing $\mathbf{S}_{\mathrm{next}}$. 
        \label{code:softmax}
        \STATE Await completion of the $\tilde{\mathbf{P}}_{\mathrm{cur}} \mathbf{V}_{j-1}$ WGMMA step, then readjust scaling of $\mathbf{O}_i$
        \STATE Release the $(j\,\%\,s)$th and $(j-1\,\%\,s)$th positions in the buffers for $\mathbf{K}$ and $\mathbf{V}$, respectively.
        \STATE Assign $\mathbf{S}_{\mathrm{cur}}$ the value of $\mathbf{S}_{\mathrm{next}}$.
        \label{code:mainloop_end}
      \ENDFOR 
      \STATE Hold until $\mathbf{V}_{T_c - 1}$ has been loaded into shared memory.
      \STATE Calculate
      $\mathbf{O}_{i} = \mathbf{O}_{i} + \tilde{\mathbf{P}}_{\mathrm{last}} \mathbf{V}_{T_c - 1}$ by WGMMA, waiting until the operation completes.

\cref{alg:flash3_wgmma} serves as the substitute for the consumer procedure in \cref{alg:flash3_wgmma_ws_only}, thereby forming the complete \textsc{FlashAttention-3} methodology in FP16 format. Conceptually, we employ WGMMA as shorthand for asynchronous GEMM computations. In the main processing loop (spanning lines \ref{code:mainloop_start} through \ref{code:mainloop_end}), the second invocation of WGMMA in the $j$-th iteration (refer to line \ref{code:second_wgmma}) is executed concurrently with the softmax calculation of the subsequent iteration $j+1$ (refer to line \ref{code:softmax}).

Although the pipelined architecture depicted above provides promising theoretical benefits, its real-world application necessitates attention to several practical challenges:
\paragraph{Compiler reordering} While the pseudocode reflects an ideal sequence of operations, in reality, the NVCC compiler often reorganizes instructions for efficiency purposes. This instruction reordering may compromise the intended pipelining between WGMMA and non-WGMMA operations, which can undermine the expected behavior or lessen the realized performance improvements. SASS code inspection confirms that the generated assembly does maintain the expected overlap (see Section~\ref{sec:2-stage-sass}).
\paragraph{Register pressure} Achieving peak performance requires the avoidance of register spilling. The 2-stage pipeline, however, introduces an elevated need for register storage to retain both intermediate results and pipeline context. In particular, an extra $\mathbf{S}_{\mathrm{next}}$ value must remain resident in registers, incurring an additional cost of $B_r \times B_c \times \text{sizeof}(\text{float})$ bytes for each threadblock. This uptick in register utilization can limit the feasibility of using larger thread block sizes, as they are also heavily dependent on available registers. Profiling should be used to judiciously select the most appropriate parameter settings.
\paragraph{3-stage pipelining} Building on the earlier 2-stage pipeline, we introduce a 3-stage alternative that permits further overlap of the second WGMMA operation with the softmax computation. This design has the potential to push Tensor Core occupancy even higher, though it comes with the downside of requiring yet more registers—heightening the tension between selecting large tile sizes and maximizing pipeline depth. Section~\cref{sec:3-stage} provides an in-depth account of the 3-stage variant along with its empirical assessment.

\subsection{Low-precision with FP8}
\label{sec:algofp8}

\textbf{Efficiency: layout transformations.} Executing the forward pass of \textsc{FlashAttention-3} using FP8 precision introduces new complications regarding layout compliance that do not arise with FP16.

Typically, input tensors $\mathbf{Q}$, $\mathbf{K}$, and $\mathbf{V}$ are arranged so that the data is contiguous along the head dimension. However, for the second GEMM within FP8 WGMMA, the k-major ordering requires that $\mathbf{V}$—more precisely, the tiles of $\mathbf{V}$ that are loaded into SMEM—be stored contiguously along the sequence length dimension. As the TMA load does not alter which dimension is contiguous, this necessitates either (1) transposing $\mathbf{V}$ in GMEM ahead of computation as a pre-processing step, or (2) performing an in-kernel tile-wise transpose on $\mathbf{V}$ after transferring tiles into SMEM. For approach (1), one can either (1a) combine the transpose with the epilogue of a preceding computation, such as after rotary embedding, or (1b) invoke a dedicated pre-processing transpose kernel\footnote{A high-performance transpose kernel can achieve throughput close to the theoretical device bandwidth~\citep{colfax_cutlass_transpose_2024}.} to swap the strides of the sequence length and head axes. Nonetheless, (1a) presents integration difficulties when working with a conventional library setup, and (1b) introduces significant overhead, making it particularly unsuitable when memory bandwidth is the bottleneck, as is often the case during inference.

In contrast, for FP8 \textsc{FlashAttention-3}, we select option (2). Within the in-kernel transpose, we utilize the LDSM (\verb|ldmatrix|) and STSM (\verb|stmatrix|) instructions. These operations allow a warp of threads to cooperatively transfer data between shared memory (SMEM) and register memory (RMEM) in chunks of 128 bytes.\footnote{The PTX manual notes that LDSM/STSM are intended for copying $8 \times 8$ matrices of 16-bit elements \cite[\S 9.7.13.4.15-16]{ptx}. Nonetheless, by packing two 8-bit values together, we can leverage these instructions with FP8 data. It is worth noting that the transpose variants of LDSM/STSM cannot break apart such packed 8-bit elements, leading to a requirement for specific register shuffles between the LDSM and STSM to achieve true tile-wise transposition. We omit these implementation specifics here.} Both LDSM and STSM exhibit efficient register usage and are suitable for execution directly within the producer warpgroup. Additionally, they can handle layout transpositions during memory copying. Furthermore, after completing the initial iteration, we schedule the transpose of the subsequent $\mathbf{V}$ tile to take place concurrently with the execution of the two WGMMAs that process the preceding $\mathbf{V}$ tile along with the current $\mathbf{K}$ tile.

Second, it is notable that, in contrast to FP16, the FP32 accumulator used within an FP8 WGMMA presents a memory layout that diverges from the expectation for operand A when held within registers. This distinction is illustrated in \cref{fig:rmem_accum} and \cref{fig:rmem_operand}, where the figures show the order in which register-resident entries are associated with each thread. Utilizing byte permutation instructions, it becomes possible to reorganize the first WGMMA's accumulator into a structure compatible both for a subsequent WGMMA and with the particular arrangement of the $\mathbf{V}$ tile emerging from the in-kernel transposition. Specifically, as shown in \cref{fig:rmem_accum}, we re-sequence the order to
$$\{ \verb|d0 d1 d4 d5 d2 d3 d6 d7| \},$$
and this rearrangement of registers is systematically applied for every 8-byte group. At the tile level for $\mathbf{P}$, such a transformation results in permuting its columns, so that, for example, columns $0189$ are relocated as the leading four columns. In order for WGMMA to produce the correct tile output, we ensure that the in-kernel transposition writes the $\mathbf{V}$ tile using a correspondingly permuted row sequence.\footnote{With this greater flexibility during the in-kernel transposition, there is no need to use shuffle instructions to transfer register data between threads, in contrast to earlier strategies described in~\citep{colfax_fp8_flashattention_2024}.}

\textbf{Accuracy: block quantization and incoherent processing.} The FP8 (e4m3) format encodes numbers using 3 bits for the mantissa and 4 bits for the exponent, leading to greater numerical inaccuracies compared to FP16 or BF16 representations. Additionally, large-scale models are known to exhibit extreme outlier values~\citep{dettmers2208llm, sun2024massive} that can vastly exceed the range of typical activations, complicating effective quantization. A common mitigation strategy is per-tensor scaling~\citep{micikevicius2022fp8}, where each tensor (such as $\mathbf{Q}$, $\mathbf{K}$, and $\mathbf{V}$) is assigned a single scaling factor. In order to enhance the accuracy of FP8-based attention computations, we make use of two specific methods:

\begin{enumerate}

\item \textbf{Block quantization}: In this approach, a single scalar value is stored for each block of the tensor. For each of $\mathbf{Q}$, $\mathbf{K}$, and $\mathbf{V}$, the tensor is partitioned into blocks sized $B_r \times d$ or $B_c \times d$, and each block is quantized independently. This quantization process can be integrated with preceding operations, such as rotary embedding, without incurring additional computational overhead, as rotary embedding is constrained by memory bandwidth rather than compute. Given that the \textsc{FlashAttention-3} algorithm is inherently block-based, it is possible to rescale each block of $\mathbf{S}$ to accommodate the effects of block quantization without increasing computational costs.

\item \textbf{Incoherent processing}: To mitigate the influence of outlier values, $\mathbf{Q}$ and $\mathbf{K}$ are pre-multiplied by a random orthogonal matrix $\mathbf{M}$ prior to FP8 quantization. Because the orthogonality of $\mathbf{M}$ ensures $\mathbf{M} \mathbf{M}^\top = I$, it follows that $(\mathbf{Q} \mathbf{M}) (\mathbf{K} \mathbf{M})^\top = \mathbf{Q} \mathbf{K}^\top$; that is, applying $\mathbf{M}$ to both $\mathbf{Q}$ and $\mathbf{K}$ leaves the result of the attention computation unchanged. This operation helps distribute outlier magnitudes more evenly since each component of $\mathbf{Q} \mathbf{M}$ or $\mathbf{K} \mathbf{M}$ is formed as a random linear combination of the original values, thereby lessening quantization-induced errors. Consistent with the approaches of \citet{chee2024quip} and~\citet{tseng2024quip}, $\mathbf{M}$ is chosen as a composition of random diagonal matrices with entries $\pm 1$ and a Hadamard matrix. This structure allows multiplication to be performed in $O(d \log d)$ time instead of $O(d^2)$, and such multiplications can also be efficiently combined with rotary embedding without added computational expense.
\end{enumerate}

We validate that these two techniques reduces numerical error by up to 2.6$\times$ in \cref{sec:numerical_error}.

\section{Empirical Validation}
\label{sec:exp}

We use the primitives from CUTLASS~\citep{Thakkar_CUTLASS_2023} such as WGMMA and TMA abstractions to implement \textsc{FlashAttention-3} and evaluate its efficiency and accuracy.

\begin{itemize}

\item \textbf{Attention performance evaluation.} We assess the execution time of \textsc{FlashAttention-3} over a range of sequence lengths, comparing its performance against a standard PyTorch version, \textsc{FlashAttention-2}, a Triton-based \textsc{FlashAttention-2} implementation leveraging H100-specific features, and a vendor-optimized H100 cuDNN variant of \textsc{FlashAttention-2}. Our experiments demonstrate that \textsc{FlashAttention-3} achieves up to a 2.0$\times$ speedup relative to \textsc{FlashAttention-2}, and a 1.5$\times$ improvement over the Triton implementation. Additionally, \textsc{FlashAttention-3} attains throughput as high as 740 TFLOPs/s, reaching 75\% of the H100 GPU’s peak theoretical throughput.

\item \textbf{Ablation analysis.} We provide evidence that both warp-specialization and GEMM-softmax pipelining within our algorithm are key contributors to the increased performance observed for \textsc{FlashAttention-3}.

\item \textbf{FP8 attention accuracy.} Our findings show that employing block quantization alongside incoherent processing effectively reduces the numerical error for FP8 \textsc{FlashAttention-3} by a factor of 2.6$\times$.

\subsection{Benchmarking Attention}
\label{subsec:benchmark_attn}

We evaluate the performance of various attention mechanisms on an H100 80GB SXM5 GPU under multiple configurations, specifically toggling between scenarios with and without causal masks and using head dimensions of either 64 or 128, while employing FP16 input formats. The findings, presented in~\cref{fig:benchmark_attn_fwd} and~\cref{fig:benchmark_attn_bwd}, indicate that \textsc{FlashAttention-3} outperforms \textsc{FlashAttention-2} by approximately 1.5-2.0$\times$ in forward computation and 1.5-1.75$\times$ in the backward direction. When compared against the traditional attention implementation, \textsc{FlashAttention-3} achieves speedups of up to 3-16$\times$. Notably, for medium and large sequence lengths (1k and beyond), \textsc{FlashAttention-3} even outpaces a proprietary, highly optimized library (cuDNN) specifically tailored for H100 GPUs.

\paragraph{Benchmark settings:} The experiments adjust sequence length values across 512, 1k, up to 16k, maintaining a batch size such that the aggregate token count remains fixed at 16k. The hidden layer dimension is fixed at 2048, with each attention head dimension chosen from 64, 128, or 256, corresponding to 32, 16, or 8 heads, respectively. FLOPs for the forward computation are determined using:

\begin{equation*}
  4 \cdot \text{seqlen}^2 \cdot \text{head dimension} \cdot \text{number of heads}.
\end{equation*}

When applying causal masking, we adjust the value by dividing by 2, reflecting that only about half of the matrix entries are actually computed. For the backward pass FLOPs, we obtain this metric by scaling the forward pass FLOPs by a factor of 2.5; this is due to the fact that the forward computation involves two matrix multiplications (matmuls), whereas the backward computation—due to recomputation—requires five matmuls.

We also evaluate the forward pass runtime using FP8 precision under comparable conditions. The results for headdim 256 are displayed in ~\cref{fig:benchmark_attn_fp8}, with comprehensive findings presented in \cref{sec:benchmark_attn_fp8_full}.

\subsection{Ablation Study: 2-Stage Pipelining Experiments}
\label{sec:2-stage-pipelining-experiments}

We conduct an ablation study of both the 2-stage WGMMA-softmax pipeline and warp-specialization in the context of non-causal FP16 \textsc{FlashAttention-3}, maintaining fixed parameters $\{\text{batch}, \text{seqlen}, \text{nheads}, \text{hdim}\} = \{ 4, 8448, 16, 128\}$. The findings, presented in~\cref{table:ablation_pipelining}, demonstrate that the proposed algorithmic enhancements—namely, introducing asynchrony via warp-specialization and enabling overlap between GEMM and softmax—produce a substantial performance gain, elevating throughput from 570 to 661 TFLOPs.

\subsection{Numerical Error Validation}
\label{sec:numerical_error}

Given the ongoing attention to the numerical accuracy~\citep{golden2024flash} of \textsc{FlashAttention}, we evaluate \textsc{FlashAttention-2}, \textsc{FlashAttention-3}, and a conventional attention implementation in comparison to a FP64 reference implementation. In order to mimic the presence of outlier features and activations typically found in LLMs~\citep{dettmers2208llm, sun2024massive}, we sample the elements of $\mathbf{Q}, \mathbf{K}, \mathbf{V}$ according to the following distribution:

\begin{equation*}
  \mathcal{N}(0, 1) + \mathcal{N}(0, 100) \cdot \mathrm{Bernoulli}(0.001).
\end{equation*}

In this setup, each matrix element follows a normal distribution with zero mean and standard deviation of 1; additionally, 0.1\% of the entries include an extra, independent noise term sampled from a normal distribution with a standard deviation of 10. We report the root mean squared error (RMSE) results in~\cref{table:numerical_error}. When using FP16 precision, both \textsc{FlashAttention-2} and \textsc{FlashAttention-3} yield RMSE values that are 1.7$\times$ lower than those from the conventional implementation, as intermediate computations (specifically the softmax outputs) are performed in FP32. The baseline FP8 attention approach applies per-tensor scaling, utilizes an FP32 accumulator during matrix multiplication, and stores intermediate softmax values in FP16. Through the use of block quantization and incoherent processing, \textsc{FlashAttention-3} operating in FP8 achieves RMSE that is 2.6$\times$ smaller than that of the baseline.

\section{Dicussion, Limitations, Conclusion}
\label{sec:discussion}

Through the development of \textsc{FlashAttention-3}, we have shown that employing modern programming methods and leveraging advanced hardware capabilities like asynchronous operations and reduced-precision arithmetic can significantly enhance both the performance and fidelity of attention mechanisms. Our approach achieves a 1.5-2.0$\times$ acceleration in attention computation relative to \textsc{FlashAttention-2}, and delivers a 2.6$\times$ decrease in FP8 numerical errors when compared to traditional per-tensor quantization techniques. There remain several aspects for improvement in future work: optimizing for large language model (LLM) inference, implementing a persistent kernel strategy within the FP8 kernel,\footnote{For our reported results, FP16 \textsc{FlashAttention-3} employs a persistent kernel alongside load balancing, whereas this feature is absent in the FP8 version. This distinction partially accounts for FP8 \textsc{FlashAttention-3}'s reduced performance for small sequence lengths and under causal masking relative to FP8 cuDNN kernels.} and investigating how reduced-precision attention affects large-scale training regimens. While this study centers on Hopper GPUs, we anticipate that our methodology will be transferable to other accelerator platforms as well. Ultimately, we hope that enhanced primitives for attention will foster the development of applications that require processing long contexts.

\end{document}